\section{GenAI 技术演进时间线}

\subsection{2018-2021: foundation model 时代}

这一阶段最大的突破是 Transformer 预训练+微调。Burak 讲到当时的 GTM：先用 BERT、GPT-2 解决 NLU/NLG 任务，再把 embedding 开放到企业 Search 与 document understanding。这个阶段的典型做法是 building an "AI co-pilot" for narrow tasks，强调「数据准备、feature engineering、knowledge graph」才是真正的竞争力。

\subsection{2022-2023: Agent 与 multi-turn}
Agent 的概念在 2022 年开始普及：LLM 不再只输出文本，而是 orchestrating tools。Vertex AI Agent Builder、LangChain 等框架出现后，工程师可以在几天内拼出 plan→tool→feedback 流程。这个阶段的痛点是「hard-coded workflow + manual monitoring」，所以 Burak 强调要做 metrics-first design。

这个阶段常见的 repository 包含：prompt templates、tool manifest、evaluation harness。Burak 特别提到：要把 tool 的 sample input/output, estimated latency, failure modes 都写成 doc 供 planner 参考。否则 agent 容易因误判 tool call 导致 cascade failure。

\subsection{2024 之后: 实时、可控、合规}

进入 2024 年，企业关注点转向实时性与合规。Gemini 的 multi-modal、Vertex 的 agent scheduler、Anthropic 的 RSP，都对「model governance」提出更高要求。新阶段的关键词是「可追溯」、可视化「policy-as-code」与严谨的「failure analysis」。

一个关键趋势是「legal hold + audit trail」逐渐成为 baseline。合规团队要求每个 decision path 都可以 reproduce，甚至要把 user prompt 与 tool call 串成 timeline 供审查。很多企业把这挂到 data lake，和 observability pipeline 联动。

\subsection{本章小结}

回顾时间线可以看到：从预训练到 agent，再到合规，企业 AI 的挑战一路从 capability 转向 control，这意味着团队必须不断迭代流程与组织结构。

